\begin{lemma}
Let $n\ge0$ be an integer, write $\mathbb{B}=\{\xi:|\xi|<1\}$, and let
$H\colon\mathbb{C}\to\mathbb{C}$ satisfy
\begin{enumerate}
\item[(H1)] $H$ is continuous on $\mathbb{B}$;
\item[(H2)] $H\in W^{1,2}_{\mathrm{loc}}(\mathbb{B})$ in the sense of
  \noderef{MemWLocOn};
\item[(H3)] $H$ is differentiable at almost every point of $\mathbb{B}$.
\end{enumerate}
Let $a_k$, $k\in\mathbb{Z}$, be the Fourier coefficients of
\noderef{FourierCoefOfH} of $H$, and write $A_k$ for the Fourier coefficient of
the same index of $H_r$, the polar radial derivative of \noderef{HRad}, and
$B_k$ for that of $H_\theta$, the polar angular derivative of
\noderef{HAng}.  Then
\begin{enumerate}
\item[(a)] each $a_k$ is continuous on $[0,1)$;
\item[(b)] for every $k$ and all $c,d\in(0,1)$, the function $A_k$ is
  integrable on the closed interval with endpoints $c$ and $d$, and
  $a_k(d)-a_k(c)=\displaystyle\int_c^dA_k(s)\,ds$;
\item[(c)] for every $k$ and almost every $r\in(0,1)$, $a_k$ is differentiable
  at $r$ with $a_k'(r)=A_k(r)$;
\item[(d)] for every $k$ and almost every $r\in(0,1)$, $B_k(r)=i(k-n)a_k(r)$;
\item[(e)] for almost every $r\in(0,1)$ the functions
  $\theta\mapsto H_r(re^{i\theta})$ and $\theta\mapsto H_\theta(re^{i\theta})$
  are integrable on $[0,2\pi]$;
\item[(f)] for every $r\in[0,1)$ the expansion (29) holds in $L^2(0,2\pi)$:
\[
  \int_0^{2\pi}\Bigl|H\bigl(re^{i\theta}\bigr)
    -\sum_{k=-N}^{N}a_k(r)e^{i(k-n)\theta}\Bigr|^2\,d\theta\ \longrightarrow\ 0
  \qquad(N\to\infty).
\]
\end{enumerate}

This is the paper's (31) and (32) with Section 3's $C^\infty$ standing
assumption removed.  The paper reads that assumption at exactly one place,
paper.tex:265--266, to obtain $a_k\in C[0,1)\cap C^\infty(0,1)$, whereas
paper.tex:269 justifies the term-by-term differentiation (32) from
$H\in W^{1,2}(\mathbb{B})$ alone.  Accordingly the $C[0,1)$ half of (31) is (a),
which needs only continuity -- free for the inverse of a quasiconformal
homeomorphism -- while the $C^\infty(0,1)$ half is deliberately \emph{not}
claimed here: it is regained downstream from the system (33) that the
coefficients satisfy, which is the paper's own reason for (31) in the first
place.  Parts (c) and (d) are what (32) asserts, namely that the Fourier
coefficients of $H_r$ and $H_\theta$ are $a_k'$ and $i(k-n)a_k$; (b) is the
integrated form of (c), which is what survives when $a_k$ is not yet known to
be differentiable everywhere; (e) is the integrability on circles that the
passage from (28) to (33) needs in order to integrate in $\theta$; and (f) is
the expansion (29).
\end{lemma}
\begin{proof}
Fix $k\in\mathbb{Z}$ and write, for $r\in[0,1)$ and $\theta\in\mathbb{R}$,
\[
  \gamma_r(\theta)=re^{i\theta},\qquad \varphi_k(\theta)=e^{i(n-k)\theta},
\]
so that, by \noderef{FourierCoefOfH},
\begin{equation}\tag{0}
  a_k(r)=\frac1{2\pi}\int_0^{2\pi}H\bigl(\gamma_r(\theta)\bigr)\varphi_k(\theta)\,d\theta .
\end{equation}
Since $|e^{i\theta}|=1$ we have $|\gamma_r(\theta)|=r<1$, so
$\gamma_r(\theta)\in\mathbb{B}$ for every $\theta$; and $|\varphi_k(\theta)|=1$.

\emph{Step 1: joint continuity and local bounds.}  The map
$(r,\theta)\mapsto\gamma_r(\theta)$ is continuous from $[0,1)\times\mathbb{R}$
into $\mathbb{B}$, so $(r,\theta)\mapsto H(\gamma_r(\theta))\varphi_k(\theta)$
is continuous on $[0,1)\times\mathbb{R}$, by (H1) and by continuity of
$\varphi_k$.  For $0\le\rho<1$ the set $\{\xi:|\xi|\le\rho\}$ is a compact
subset of $\mathbb{B}$, so
\[
  M_\rho=\sup\bigl\{|H(\xi)|:|\xi|\le\rho\bigr\}<\infty,
\]
and $|H(\gamma_r(\theta))\varphi_k(\theta)|\le M_\rho$ for all $r\in[0,\rho]$
and all $\theta$.

\emph{Step 2: proof of (a).}  Let $r_0\in[0,1)$ and choose $\rho$ with
$r_0<\rho<1$.  For $r\in[0,\rho]$ the integrand in (0) is dominated by the
constant $M_\rho$, which is integrable on $[0,2\pi]$, and for each fixed
$\theta$ it is continuous in $r$ by Step 1.  Hence for any sequence
$r_j\to r_0$ in $[0,\rho]$ the dominated convergence theorem gives
\[
  \int_0^{2\pi}H\bigl(\gamma_{r_j}(\theta)\bigr)\varphi_k(\theta)\,d\theta
  \longrightarrow
  \int_0^{2\pi}H\bigl(\gamma_{r_0}(\theta)\bigr)\varphi_k(\theta)\,d\theta,
\]
so $a_k$ is continuous at $r_0$.  As $r_0\in[0,1)$ was arbitrary, $a_k$ is
continuous on $[0,1)$, and $k$ was arbitrary; this is (a).

\emph{Step 3: the punctured disk as an annulus, and the Sobolev data there.}
Put
\[
  A=\{\xi:0<|\xi|<1\}=\mathbb{B}-\{0\},
\]
the annulus with radii $a=0$ and $b=1$, so that $0\le a<b$.  By (H2) and
\noderef{MemWLocOn} there is a family $D$ of functions indexed by words in
$\{0,1\}$ with $D[\,]=H$, with $D(w)\in L^{2}_{\mathrm{loc}}(\mathbb{B})$ in the
sense of \noderef{MemLpLocOn} for every word $w$ of length at most $1$, and
with $D([j])$ the distributional $j$-th partial derivative of $H$ on
$\mathbb{B}$ in the sense of \noderef{IsWeakPartialOn} for $j=0,1$ --- the
latter because the only word $w$ with $|w|<1$ is the empty word, and
$[\,]+\!\!+[j]=[j]$.  Put $D_j=D([j])$.  We check the four hypotheses of
\noderef{CircleAvgLocAbsCont} for $n$, $a=0$, $b=1$, the index $k$, the
function $H$ and the family $D_0,D_1$.

For its (H1): let $\kappa\subseteq A$ be compact.  Then $\kappa$ is a compact
subset of $\mathbb{B}$, so by the above $\mathbf1_\kappa H$ and
$\mathbf1_\kappa D_j$ lie in $L^{2}$; and $\kappa$ has finite Lebesgue measure,
so by the Cauchy--Schwarz inequality $\int_\kappa|g|\le
\bigl(\int_\kappa|g|^{2}\bigr)^{1/2}\sigma(\kappa)^{1/2}<\infty$ for
$g\in\{H,D_0,D_1\}$.  Hence $H$ and the $D_j$ lie in $L^{1}_{\mathrm{loc}}(A)$
in the sense of \noderef{MemLpLocOn}.

For its (H2): a test function for $A$ is a test function for $\mathbb{B}$,
since by \noderef{IsTestFn} the only clause mentioning the set is
$\operatorname{tsupp}\varphi\subseteq A\subseteq\mathbb{B}$, and the
integration-by-parts identity of \noderef{IsWeakPartialOn} does not mention the
set; so $D_j$ is the distributional $j$-th partial derivative of $H$ on $A$ as
well.

For its (H3): $H$ is continuous on $A\subseteq\mathbb{B}$ by (H1) of the
present lemma.

For its (H4): by (H3) of the present lemma $H$ is differentiable at almost
every point of $\mathbb{B}$, hence at almost every point of the subset $A$.

\emph{Step 4: proof of (b), (c), (d), (e).}  \noderef{CircleAvgLocAbsCont},
applied as in Step 3, yields for this $k$: its (a), that $A_k$ is integrable on
the closed interval with endpoints $c$ and $d$ for all $c,d\in(0,1)$; its (b),
that $a_k(d)-a_k(c)=\int_c^dA_k(s)\,ds$ for all such $c,d$; its (c), that for
almost every $r\in(0,1)$ the function $a_k$ is differentiable at $r$ with
$a_k'(r)=A_k(r)$; its (d), that $B_k(r)=i(k-n)a_k(r)$ for almost every
$r\in(0,1)$; and its (e), that for almost every $r\in(0,1)$ the functions
$\theta\mapsto H_r(re^{i\theta})$ and $\theta\mapsto H_\theta(re^{i\theta})$ are
integrable on $[0,2\pi]$.  The first two are (b), the third is (c), the fourth
is (d), all for this $k$ and hence, $k$ being arbitrary, for every $k$; and the
fifth is (e), which does not depend on $k$.

\emph{Step 5: proof of (f).}  Fix $r\in[0,1)$.  By Step 1 the function
$\theta\mapsto H(\gamma_r(\theta))$ is continuous on the compact interval
$[0,2\pi]$, hence bounded and measurable there, so it belongs to
$L^2(0,2\pi)$.  The family $(e_m)_{m\in\mathbb{Z}}$,
$e_m(\theta)=(2\pi)^{-1/2}e^{im\theta}$, is an orthonormal basis of
$L^2(0,2\pi)$.  Its $m$-th coefficient for our function is
\[
  \bigl\langle H\circ\gamma_r,\,e_m\bigr\rangle
  =(2\pi)^{-1/2}\int_0^{2\pi}H\bigl(\gamma_r(\theta)\bigr)e^{-im\theta}\,d\theta .
\]
Taking $m=k-n$ and using $e^{-i(k-n)\theta}=e^{i(n-k)\theta}=\varphi_k(\theta)$,
this equals $(2\pi)^{1/2}a_k(r)$ by (0); consequently
\[
  \bigl\langle H\circ\gamma_r,\,e_{k-n}\bigr\rangle\,e_{k-n}(\theta)
  =(2\pi)^{1/2}a_k(r)\cdot(2\pi)^{-1/2}e^{i(k-n)\theta}
  =a_k(r)\,e^{i(k-n)\theta}.
\]
Because $(e_m)_{m\in\mathbb{Z}}$ is an orthonormal basis, the expansion of
$H\circ\gamma_r$ converges unconditionally in $L^2(0,2\pi)$: the net of partial
sums over finite subsets of the index set converges to $H\circ\gamma_r$.  In
particular, along any sequence of finite index sets increasing to the whole
index set the partial sums converge in $L^2(0,2\pi)$.  The sets
$\{k\in\mathbb{Z}:|k|\le N\}$ increase to $\mathbb{Z}$ as $N\to\infty$, and the
corresponding partial sum is, by the display above,
$\sum_{k=-N}^{N}a_k(r)e^{i(k-n)\theta}$.  Hence
\[
  \int_0^{2\pi}\Bigl|H\bigl(\gamma_r(\theta)\bigr)
    -\sum_{k=-N}^{N}a_k(r)e^{i(k-n)\theta}\Bigr|^2\,d\theta\longrightarrow0
\]
as $N\to\infty$, which is (f).
\end{proof}
